\par\medskip
\noindent
\underline{Teorema} 5.2. \underline{Supongamos que $S$ es un rasgo
completo, y sean $A_K$ un esquema abeliano sobre $K$, $A'_K$ el esquema
abeliano dual y $\ell$ un primo tal que el modelo de Néron conexo $A^o$
es $\ell$-divisible (cf. 2.2.1), de modo que se tienen entonces las
filtraciones (5.1.10) y (5.1.11) de los progrupos $T_\ell(A_K)$ y
$T_\ell(\sisi{A_{K'}}{A'_K})$. Entonces las partes tóricas se recuperan de las partes
fijas mediante la fórmula}
\begin{equation}\tag{5.2.1}
 T_\ell(A_K)^t
 =
 T_\ell(A_K)^f\cap\bigl(T_\ell(A'_K)^f\bigr)^\perp,
\end{equation}
\underline{y mediante la fórmula análoga obtenida intercambiando los
papeles de $A_K$ y $A'_K$, donde $\perp$ denota el complemento ortogonal
respecto del emparejamiento canónico (1.0.2).}

\underline{Por último, cuando $A_K$ posee reducción semiestable sobre
$S$ (lo que se deduce de las demás hipótesis en el caso $\ell=p$), se
tiene incluso la fórmula}
\begin{equation}\tag{5.2.2}
 T_\ell(A_K)^t=\bigl(T_\ell(A'_K)^f\bigr)^\perp.
\end{equation}

\smallskip
\noindent
5.2.3. El significado de la fórmula (5.2.1) es muy claro cuando $\ell$
es primo con la característica de $K$: como en 2.4, puede interpretarse
entonces pasando a los $\mathbb Z_\ell$-módulos ordinarios, dotados de
una acción de $\pi=\operatorname{Gal}(\overline K/K)$, formados por los
puntos con valores en $\overline K$ de los progrupos que intervienen.
Expliquémoslo en el caso general, utilizando la noción de dual
\begin{equation}\tag{5.2.3.1}
 D(G)=\operatorname{Hom}\bigl(G,T_\ell(\mathbb G_m)\bigr)
\end{equation}
de un pro-$\ell$-grupo de Barsotti--Tate $G=(G(\nu))$. Por definición,
es el sistema proyectivo de los duales de Cartier $D(G(\nu))$
(SGA 3 VII 3.3.1 y VIII 1), cuyos morfismos de transición se definen
de la manera evidente (por ejemplo, por transposición a partir de los
morfismos de inclusión entre los $G(\nu)$). La teoría de dualidad de
las variedades abelianas nos dice que el emparejamiento mencionado en
5.2 define, de hecho, un \underline{isomorfismo}
\begin{equation}\tag{5.2.3.2}
 T_\ell(A'_K)\simeq D\bigl(T_\ell(A_K)\bigr),
\end{equation}
y por tanto, por composición, un morfismo
\begin{equation}\tag{5.2.3.3}
 T_\ell(A_K)^f\longrightarrow D\bigl(T_\ell(A'_K)^f\bigr).
\end{equation}
La fórmula (5.2.1) expresa que el núcleo del homomorfismo precedente
(que a priori es un prosubgrupo de su primer término) es precisamente
la parte tórica $T_\ell(A_K)^t$. Obsérvese que se trata de un
\underline{isomorfismo de proobjetos, y no de sistemas proyectivos}:
no afirmamos (y en general sería falso) que, para todo entero $\nu$,
$T_\ell(A_K)^t(\nu)$ sea igual al núcleo de
\[
 T_\ell(A_K)^f(\nu)
 \longrightarrow
 D\bigl(T_\ell(A'_K)^f\bigr)(\nu);
\]
este último puede ser a priori mayor. Veremos más adelante, sin embargo,
que no es mayor si $A_K$ posee reducción semiestable; por las hipótesis
formuladas en 5.2, este es, en particular, el caso cuando $\ell=p$.

\smallskip
\noindent
5.2.4. La demostración de (5.2.1) no hace sino repetir, casi sin cambios,
la demostración de 2.4. Obsérvese que el emparejamiento inducido
\begin{equation}\tag{5.2.4.1}
 T_\ell(A_K)^f\times T_\ell(A'_K)^f
 \longrightarrow T_\ell(\mathbb G_{mK})
\end{equation}
que da lugar a (5.2.3.3) está inducido por un emparejamiento
\begin{equation}\tag{5.2.4.2}
 T_\ell(A^o)^f\times T_\ell({A'}^o)^f
 \longrightarrow T_\ell(\mathbb G_{mS})
\end{equation}
de progrupos de Barsotti--Tate sobre $S$, inducido a su vez por
el emparejamiento de progrupos
\begin{equation}\tag{5.2.4.3}
 T_\ell(A^o)\times T_\ell({A'}^o)
 \longrightarrow T_\ell(\mathbb G_{mS})
\end{equation}
deducido de la biextensión $W$ de $(A^o,{A'}^o)$ por
$\mathbb G_{mS}$. En consecuencia, el emparejamiento
\begin{equation}\tag{5.2.4.4}
 T_\ell(A_o^o)\times T_\ell({A'}_o^o)
 \longrightarrow T_\ell(\mathbb G_{mk}),
\end{equation}
obtenido de (5.2.4.2) al pasar a las fibras especiales, es precisamente
el emparejamiento deducido de la biextensión $W_o$ de
$(A_o^o,{A'}_o^o)$ por $\mathbb G_{mk}$. Por VIII 4.11 y 1.5, el
anulador izquierdo del emparejamiento (5.2.4.4) es, por tanto,
$T_\ell(T_o)$, que es precisamente la fibra especial de la parte tórica
$T_\ell(A^o)^t$. Volviendo a la definición de anuladores izquierdos
dada en 5.2.3, concluimos mediante el lema siguiente.

\medskip
\noindent
\underline{Lema} 5.2.5. \underline{Sea $u:G\longrightarrow H$ un
homomorfismo de pro}-$\ell$-\underline{grupos de Barsotti--Tate sobre $S$. Si
$u_o:G_o\longrightarrow H_o$ es cero, entonces $u$ es cero. Si
$u_o:G_o\longrightarrow H_o$ es una isogenia (respectivamente, un
isomorfismo), entonces $u$ es una isogenia (respectivamente, un
isomorfismo).}

En efecto, apliquemos la primera afirmación de 5.2.5 al morfismo
compuesto
\[
 T_\ell(A^o)^t
 \longrightarrow T_\ell(A^o)^f
 \longrightarrow D\bigl(T_\ell({A'}^o)^f\bigr),
\]
donde el segundo morfismo se deduce del emparejamiento (5.2.4.2). El
compuesto es cero, de modo que el emparejamiento considerado se anula
sobre la parte tórica de la izquierda. Por simetría, se anula también
sobre la parte tórica de la derecha y desciende, por tanto, a un
emparejamiento sobre los cocientes, o \emph{trozos abelianos},
\[
 T_\ell(A^o)^{ab}\times T_\ell({A'}^o)^{ab}
 \longrightarrow T_\ell(\mathbb G_m).
\]
Estos cocientes son de nuevo progrupos de Barsotti--Tate, como
se comprueba fácilmente. Queda demostrar que el morfismo correspondiente
a este emparejamiento,
\[
 T_\ell(A^o)^{ab}
 \longrightarrow D\bigl(T_\ell({A'}^o)^{ab}\bigr),
\]
es una isogenia, lo cual se deduce nuevamente de 5.2.5.

Dejamos al lector la fácil demostración de 5.2.5, resultado válido sobre
toda base noetheriana conexa. Además, más adelante, en 7.4.7, mediante
un análisis algo más preciso de las biextensiones $W_n$ de
$(A_n^o,{A'}_n^o)$, encontraremos una demostración directa bastante
evidente de los dos casos particulares en los que tuvimos que utilizar
5.2.5.

\smallskip
\noindent
5.2.6. Supongamos ahora que $A_K$ posee reducción semiestable sobre $S$,
de modo que $\ell$ puede ser un primo cualquiera (igual o distinto de
$p$). Además, por el criterio galoisiano de reducción semiestable 3.5,
para $\ell\neq p$ se tiene
\begin{equation}\tag{5.2.7}
 \bigl(T_\ell(A'_K)^f\bigr)^\perp
 \ \bigl(=T_\ell(A_K)^t\bigr)
 \subset T_\ell(A_K)^f.
\end{equation}
Para concluir la demostración de 5.2, queda mostrar que esta inclusión
es válida sin restricción sobre $\ell$.

En efecto, el término izquierdo es canónicamente isomorfo a
\[
 D\bigl(T_\ell(A'_K)/T_\ell(A'_K)^f\bigr)
\]
y posee, por tanto, \emph{rango} (o \emph{altura})
$2n-2\alpha-\mu=\mu$, con la notación de (2.1.9), pues por hipótesis
$\lambda=0$. Este es también el rango de $T_\ell(A_K)^t$. Como la
inclusión
\begin{equation}\tag{*}
 T_\ell(A_K)^t\subset T_\ell(A_K)
\end{equation}
es \underline{estricta}, en el sentido de que ya induce un monomorfismo
sobre cada término de los sistemas proyectivos considerados, se deduce
fácilmente que el término izquierdo de (*) es maximal entre los
prosubgrupos de Barsotti--Tate de $T_\ell(A_K)$ que poseen el
mismo rango. Hay, pues, igualdad en (5.2.7); por el teorema de
ortogonalidad, esto equivale a la inclusión entre los dos términos
extremos de (5.2.7).
